\section{Sales geography is not a map of costs or assets}
\label{sec:apple-segments}

The liquidity question leaves a related operating question. If demand weakens
in one place, which part of Apple's sales and profit is exposed? Apple reports
five geographic segments, but a segment is a management reporting division,
not a list of incorporated companies, factories or product lines. Reading its
two segment notes as though those were the same partitions would make a
regional scenario appear more precise than the filing permits.

\subsection{What the regional labels measure}

Apple's fiscal 2024 Note 13 names Americas, Europe, Greater China, Japan and
Rest of Asia Pacific. Americas covers North and South America. Europe also
includes India, the Middle East and Africa; Greater China includes China
mainland, Hong Kong and Taiwan. Apple says regional net sales are
\emph{generally} based on customer location and sales through retail stores
in those geographic locations. Its segments sell similar products and
services, but are managed separately around customers, distribution partners
and market conditions \citep[Note 13, printed pp.~46--47]{apple2024}. An
India-customer sales shock therefore belongs with Apple's Europe segment
under this disclosure's general attribution rule, even though India is in Asia. Neither this mapping
nor the label ``Greater China'' locates the legal seller, the manufacturing
partner or the cash-receiving entity.

The prior revenue reading established that product and regional net sales are
two classifications of the same 391,035 million dollars in fiscal 2024 and
124,300 million in the first quarter of fiscal 2025; the totals cannot be
added. The full Note 13 and interim Note 10 tables confirm five reported
regional sales and operating-income pairs in each dated column. For example,
Greater China has fiscal 2024 sales of 66,952 million and segment operating
income of 27,082 million; in the quarter ended 28 December 2024 the amounts
are 18,513 and 8,160 million
\citep[Note 13, printed p.~47]{apple2024}
\citep[Note 10, printed p.~11]{apple2025q1}. The quarter is a distinct
three-month period, not a fourth of the fiscal-year observation. Apple's
fiscal calendar can contain 53 weeks, as it did in fiscal 2023, and its
interim MD\&A describes historically higher first-quarter sales associated
in part with the holidays \citep[Note 1 and MD\&A]{apple2024,apple2025q1}.

Management attributes the annual Greater China sales decline mainly to lower
iPhone and iPad sales and reports an unfavourable renminbi effect. For the
first quarter of fiscal 2025 it attributes the decline mainly to iPhone,
while describing a favourable currency effect
\citep[annual and interim segment operating performance]{apple2024,apple2025q1}.
These are dated issuer explanations, not measured contributions of each driver
or an estimated response equation. A future region-specific demand or
exchange-rate path needs product, currency and contract information that the
regional totals alone do not supply.

\subsection{Why segment income cannot be treated as fully allocated profit}

Apple defines each segment's operating income as third-party sales less the
related cost of sales and directly attributable operating expenses. The
annual total of segment operating income is 162,044 million, but the
consolidated operating statement reports 123,216 million. Note 13 supplies
the bridge: 31,370 million of research and development and 7,458 million of
other corporate expenses, net, are subtracted. The latter category includes
general and administrative compensation, various nonrecurring charges and
other separately managed costs. The reported relation is
\begin{equation}
162{,}044-31{,}370-7{,}458=123{,}216
\quad\text{million dollars, fiscal 2024}.
\label{eq:apple-segment-income-bridge}
\end{equation}
The first-quarter fiscal 2025 bridge has the same form:
$53{,}174-8{,}268-2{,}074=42{,}832$ million. Both endpoints agree with the
consolidated statements of operations
\citep[statements of operations and Note 13, printed pp.~29 and 47]{apple2024}
\citep[statement of operations and Note 10, printed pp.~1 and 11]{apple2025q1}.
The annual ``other corporate'' amount is not identical to the 26,097 million
selling, general and administrative expense line in the consolidated
statement. Equating the two would change the reported classification.

Equation~\eqref{eq:apple-segment-income-bridge} identifies the company-wide cost pool, not how it would be
attributed to regions. To see the difference, take a hypothetical allocation
of the 38,828 million total corporate bridge across Americas, Europe, Greater
China, Japan and Rest of Asia Pacific of respectively 14,000, 10,000, 5,000,
4,000 and 5,828 million. Moving 1,000 from Americas to Greater China produces
a second allocation of 13,000, 10,000, 6,000, 4,000 and 5,828 million.
Both are nonnegative, sum to the same corporate total and reproduce the
consolidated 123,216 million. Yet Greater China's hypothetical profit after
these costs is either $27{,}082-5{,}000=22{,}082$ or
$27{,}082-6{,}000=21{,}082$ million. Apple does not report either fully
allocated figure, and the filing supplies no rule preferring one allocation.
These constructions establish numerical nonidentification within the stated
nonnegative illustration; they do not assert that all real cost allocations
are nonnegative or that either construction describes Apple's internal books.
Forecasting regional profit also requires the behaviour of directly
attributable and corporate costs under the scenario.

\subsection{Country sales and long-lived assets answer different questions}

The annual note separately presents country sales for fiscal 2024, 2023 and
2022 and country long-lived-asset stocks at the ends of fiscal 2024 and 2023.
It separately reports countries individually accounting for at least 10\% of
the respective totals and groups the rest; that selection rule does not make
``Other countries'' a single location. The 2024 displayed amounts illustrate
the different populations:
\begin{center}
\begin{tabular}{lrr}
\toprule
Country class & 2024 net sales & Long-lived assets\\
& \multicolumn{2}{c}{USD millions}\\\midrule
U.S. & 142,196 & 35,664\\
China, including Hong Kong and Taiwan & 66,952 & 4,797\\
Other countries & 181,887 & 5,219\\\midrule
Total & 391,035 & 45,680\\\bottomrule
\end{tabular}
\end{center}
The China country-sales row equals Greater China segment sales in all three
annual columns. The note's China footnote includes Hong Kong and Taiwan, so
the equality is a useful check on these displayed populations. It does not
turn that row into a legal-entity or production-location ledger. The U.S.
country-sales row is below Americas segment sales by 24,849 million in 2024;
assigning that difference to particular countries requires the precise
country-to-region attribution and customer facts. Europe's segment label
cannot be substituted for a European-country total because it includes India,
the Middle East and Africa \citep[Note 13, printed pp.~46--47]{apple2024}.

The 45,680 million long-lived-asset total equals net property, plant and
equipment at the 2024 balance date, and the 2023 totals both equal 43,715
million at their displayed precision
\citep[Notes 5 and 13, printed pp.~39 and 47]{apple2024}. The equality is
a checked comparison, not a definition of every asset's nature, owner or
productive capacity. Country asset stocks are measured at an instant, while
sales cover a year. Further, Apple says the information supplied to its chief
operating decision maker for segment assessment excludes asset information
\citep[Note 13, printed p.~46]{apple2024}. Thus neither a segment return on
assets nor an allocation of production capacity to the five sales regions is
identified by the disclosed tables. The statement that hardware is made
substantially by outsourcing partners located primarily in China mainland,
India, Japan, South Korea, Taiwan and Vietnam concerns another population,
and does not repair that missing asset/segment mapping
\citep[Note 12, printed p.~46]{apple2024}.

\subsection{The same totals permit different product exposure}

Note 2 reports 201,183 million of fiscal 2024 iPhone sales. Its four other
product and service rows sum to 189,852 million. Along with Greater China's
66,952 million regional total, these are \emph{margins} of an unreported
product-by-region table. Let $z$ be iPhone sales assigned to Greater China,
with all amounts in millions of dollars at the displayed centres. The four
cells for Greater China versus the other four regions must be
\begin{equation}
\begin{array}{c|cc}
 &\text{iPhone}&\text{Other products and services}\\\hline
\text{Greater China}&z&66{,}952-z\\
\text{Other four regions}&201{,}183-z&122{,}900+z
\end{array}
\label{eq:apple-product-region-margins}
\end{equation}
The row and column sums in~\eqref{eq:apple-product-region-margins} hold for
any nonnegative cells; they do not select
$z$. For example, $z=36{,}000$ and $z=38{,}000$ both meet the displayed
numeric margins. For an explicit extension, let $S_r$ be each other region's
reported sales; those four regions sum to $391{,}035-66{,}952=324{,}083$.
Assign each $(201{,}183-z)S_r/324{,}083$ of iPhone sales;
these amounts sum to the remaining iPhone total and are below $S_r$.
Then split every region's remaining sales across the four non-iPhone categories
in their disclosed company-wide proportions. This produces all five
product/service and five region margins exactly. These proportional splits
are mathematical witness constructions, not issuer allocations. The source adds that the product mix
was \emph{generally} consistent across segments, with a \emph{moderately}
higher iPhone share in Greater China. Both numerical examples put Greater
China's iPhone share above the combined other-region share. The qualitative
words have no disclosed tolerance, so neither matrix is certified as a
complete interpretation of that statement
\citep[Note 2, printed pp.~35--36]{apple2024}. The interim note repeats the
qualitative comparison for its own dated quarter
\citep[Note 2, printed p.~6]{apple2025q1}; it does not supply the missing
cross-classified cells.

If a scenario \emph{assumes} a 10\% decline in Greater China iPhone sales
alone, the two illustrative fiscal 2024 matrices imply 3,600 and 3,800
million less sales, respectively. The scenario percentage is a teaching
assumption, not an estimated shock; neither sales change determines profit,
cash flow, currency translation or future behaviour. The gap between those
two conditional answers comes from the unobserved joint allocation, even
before cost responses are considered. A reproducible model should retain
the reported marginals, the qualitative mix statement, any admitted
cross-classification and its source, and the unresolved allocation when
the evidence does not determine it.

Finally, the annual and first-quarter MD\&A state Apple's then-planned
fourth-quarter fiscal 2025 adoption of ASU 2023-07 with retrospective
transition \citep[Recent Accounting Pronouncements]{apple2024,apple2025q1}.
This is a dated issuer statement. It neither supplies an already restated
2024 observation nor establishes the standard's operative provisions without
inspection of that authority. A historical query must keep the older
disclosure view separate from any later restated one when such a filing is
actually acquired. The regional calculations here are checked source-bound
comparisons and conditional witnesses, not a reconstructed Apple segment
ledger or a forecast-ready regional model.
